\documentclass[12pt,a4paper]{ctexart}
% 公共排版文件（使用 XeLaTeX 编译）
\usepackage[a4paper,margin=2.35cm]{geometry}
\usepackage{booktabs,longtable,tabularx,array,multirow,enumitem}
\usepackage{graphicx,xcolor,amsmath,hyperref,lastpage,fancyhdr}
\usepackage{tikz}
\usetikzlibrary{arrows.meta,positioning,shapes.geometric,fit,calc}
\hypersetup{colorlinks=true,linkcolor=blue!60!black,urlcolor=blue!60!black,citecolor=blue!60!black,pdftitle={逸仙活动云课程报告}}
\setlist{itemsep=2pt,topsep=3pt,leftmargin=2em}
\setlength{\parindent}{2em}
\setlength{\headheight}{15pt}
\renewcommand{\arraystretch}{1.35}
\newcolumntype{Y}{>{\raggedright\arraybackslash}X}
\newcommand{\reportcover}[1]{%
  \begin{titlepage}
  \begin{center}
  \vspace*{3cm}{\zihao{1}\bfseries 中山大学\\[0.6cm]《软件工程》期末大作业\par}
  \vspace{2.5cm}{\zihao{2}\bfseries #1\par}
  \end{center}
  \vspace{2cm}
  \begin{center}
  \begin{tabular}{l}
  {\Large 项目名称：逸仙活动云——中山大学校园活动云平台}\\
  {\Large 小组成员：钟梓俊\quad 邱剑波\quad 张宸睿}\\
  {\Large 指导老师：王可泽}
  \end{tabular}
  \end{center}
  \vfill
  \begin{center}
  {\large 2026 年 7 月 16 日\par}
  \end{center}
  \end{titlepage}}

\pagestyle{fancy}\fancyhf{}\lhead{逸仙活动云}\rhead{软件工程期末大作业}\cfoot{第 \thepage\ 页，共 \pageref{LastPage} 页}
\setcounter{tocdepth}{2}

\begin{document}
\reportcover{软件测试与质量保证报告}
\pagenumbering{Roman}\tableofcontents\clearpage\pagenumbering{arabic}

\section{测试目标}
测试的目标是验证逸仙活动云系统在正常、边界和异常条件下满足需求分析中定义的功能和非功能需求，并通过自动化回归降低迭代过程中的缺陷逃逸风险。具体包括：
\begin{enumerate}
\item 验证核心业务功能（认证、活动发现、发布审核、个人参与）的正确性和可用性。
\item 验证系统在异常输入、网络故障和越权访问下的防护能力和错误处理正确性。
\item 验证前后端接口契约的一致性，即前端期望的数据结构、状态码和错误语义与后端实现保持一致。
\item 通过自动化测试建立回归防护网，确保功能迭代不破坏已有正确行为。
\end{enumerate}

\section{测试范围}
\subsection{功能范围}
本次测试覆盖需求分析中定义的八项功能域：
\begin{enumerate}
\item \textbf{活动发现与检索（FR-01）：} 活动列表展示、关键词搜索、分类/状态筛选、分页、活动详情展示。
\item \textbf{身份认证与授权（FR-02）：} 图形验证码登录、邮箱验证码注册、JWT 会话管理、角色鉴权。
\item \textbf{活动发布与审核（FR-03）：} 草稿创建与编辑、提交审核、审核通过/驳回、状态流转。
\item \textbf{个人参与管理（FR-04）：} 报名、取消报名、收藏、取消收藏、个人日历、ICS 导出。
\item \textbf{发布者申请管理（FR-05）：} 提交申请、管理员审批、角色变更。
\item \textbf{知识关联与订阅通知（FR-06）：} 知识节点关联、订阅管理、通知生成与已读标记。
\item \textbf{数据源与内容质量（FR-07）：} 数据源维护、抓取日志、质量分、疑似重复标识、URL 安全校验。
\item \textbf{审计与安全治理（FR-08）：} 审计日志记录与查询、密码哈希存储、敏感日志脱敏。
\end{enumerate}

\subsection{非功能范围}
\begin{enumerate}
\item \textbf{性能：} 读接口 P95 响应时间不超过 500ms，写接口不超过 1s（本地/CI 环境基准）。
\item \textbf{易用性：} 页面在 360px 以上视口无横向溢出，表单校验在输入位置即时提示。
\item \textbf{安全性：} 越权请求返回 403，XSS 注入被转义，外部 URL 抓取过滤内网地址。
\end{enumerate}

\subsection{排除范围}
以下内容不纳入本次系统测试范围：第三方邮件服务的送达率和延迟、外部搜索引擎的索引质量和排序算法、AI 服务的生成内容准确性，以及数据库在高并发下的性能压力测试。这些依赖以 Mock/降级合同验证为主，不以真实第三方服务的稳定性评价本系统的正确性。

\section{测试策略}
\begin{center}
\begin{tikzpicture}[font=\small,
  layer/.style={draw,rounded corners,minimum width=9cm,minimum height=.7cm,align=center},
  arr/.style={-Latex,thick}
]
\node[layer,fill=green!10] (unit) at (0,0) {单元测试：纯函数、服务规则、组件、状态管理};
\node[layer,fill=blue!10] (integ) at (0,1.5) {集成/API 测试：Flask 测试客户端、内存 SQLite、JWT 与数据约束};
\node[layer,fill=orange!12] (e2e) at (0,3.0) {系统测试（E2E）：Playwright + Mock/Vite，验证关键用户旅程};
\node[layer,fill=red!8] (accept) at (0,4.5) {验收测试：响应式、可访问性、跨端联调、部署冒烟};
\draw[arr] (unit)--(integ);
\draw[arr] (integ)--(e2e);
\draw[arr] (e2e)--(accept);
\end{tikzpicture}
\vspace{0.2em}
{\par\small\itshape 图1：测试金字塔——从单元到验收逐层覆盖}
\end{center}

\subsection{单元测试}
前端使用 Vitest + \texttt{@vue/test-utils} 对以下模块执行单元测试：认证 Store 的登录/登出状态管理、活动排序工具函数、活动质量评分函数、日历卡片组件渲染、活动详情组件状态分支。后端使用 pytest 对服务层纯函数执行单元测试，包括内容质量评分、去重判定、安全 URL 校验、知识节点匹配和搜索日志格式化。单元测试不依赖数据库、网络和外部服务。

\subsection{集成测试}
后端集成测试使用 Flask 测试客户端和内存 SQLite 数据库，\texttt{conftest.py} 提供隔离的测试应用实例、测试 JWT、以及典型用户、活动和数据源夹具。覆盖的 API 模块包括：auth（登录、注册、用户信息）、posters（活动 CRUD、状态转换）、audit logs（审计记录写入和查询）、data sources（创建、校验、列表）、knowledge（节点管理、关联）、search（内部搜索、外部搜索降级）、quality（质量评分）、dedup（疑似重复检测）、security（URL 校验、HTML 转义）。

\subsection{系统测试}
系统测试（端到端测试）使用 Playwright 驱动 Chromium 浏览器，在 Mock 服务模式下验证完整用户旅程：访客浏览活动列表和详情页、用户注册和登录流程、发布者创建活动和提交审核、管理员审核活动和审批申请。系统测试覆盖页面跳转、权限导航、加载态和错误态等用户可见行为。

\subsection{验收测试}
验收测试以手工方式执行，覆盖三类场景：
\begin{enumerate}
\item \textbf{响应式布局：} 在 375px、768px、1024px 和桌面宽度下检查页面布局，确认无水平溢出、按钮遮挡或关键信息丢失。
\item \textbf{可访问性：} 通过 Tab 键遍历页面主要交互元素，确认焦点可见、操作可触达。
\item \textbf{部署冒烟：} 使用 Docker Compose 启动完整服务后，访问健康检查接口、执行登录和活动查询等最小业务流程。
\end{enumerate}
第三方邮件、外部搜索和 AI 服务以 Mock 验证为主，不以真实第三方稳定性评价系统正确性。

\section{测试过程}
\subsection{测试环境}
\noindent\begin{tabularx}{\textwidth}{p{2.6cm}Y}
\toprule 项目 & 规定 \\\midrule
前端环境 & Node 22（CI）/ Node 18+（本地）；\texttt{npm ci} 安装依赖；Vite 开发服务器；Mock 服务（\texttt{node mock/index.js}）；Chromium 浏览器（Playwright）。 \\
后端环境 & Python 3.10+ 虚拟环境；依赖从 \texttt{requirements-dev.txt} 安装；测试使用 \texttt{sqlite:///:memory:} 内存数据库；禁用真实 Redis、邮件服务和 LLM 调用。 \\
CI 环境 & GitHub Actions Ubuntu 最新版；\texttt{actions/setup-node@v4} 配置 Node 22；\texttt{npx playwright install --with-deps chromium} 安装浏览器。 \\
准入条件 & 相关需求编号和接口契约已确认；测试数据和测试账号可用；项目构建成功；外部依赖可 Mock 或可降级。 \\
退出条件 & P0/P1 缺陷为零；受影响自动化测试全部通过；关键手工验收用例通过；遗留风险已记录并指定负责人。 \\
\bottomrule
\end{tabularx}

\subsection{测试计划}
测试活动在项目开发周期内按迭代分阶段执行：
\begin{enumerate}
\item \textbf{第一阶段（第 1--3 周）：} 搭建测试基础设施，配置 Vitest、Playwright 和 pytest；编写认证模块（FR-02）的 API 集成测试和前端组件单元测试。
\item \textbf{第二阶段（第 4--7 周）：} 覆盖活动发布与审核（FR-03）和个人参与（FR-04）的测试用例；添加端到端测试覆盖访客浏览和登录注册流程。
\item \textbf{第三阶段（第 8--10 周）：} 补充审计（FR-08）、知识通知（FR-06）和数据治理（FR-07）的测试覆盖；执行兼容性和响应式验收测试。
\item \textbf{第四阶段（第 11--12 周）：} 全量回归测试；缺陷修复验证；测试报告和覆盖率数据归档。
\end{enumerate}

\subsection{测试用例设计}
\subsubsection{功能测试用例}
\begin{longtable}{p{1.0cm}p{2.8cm}p{4.5cm}p{4.5cm}}
\toprule 编号 & 需求/场景 & 步骤与输入 & 预期结果 \\\midrule\endhead
TC-01 & 认证成功/失败（FR-02） & 使用正确凭据登录；使用错误密码登录；缺少必填字段登录。 & 成功返回 token 和用户信息；错误密码返回 401；缺字段返回 400。 \\
TC-02 & 未认证与越权控制（FR-02） & 无 token 访问 \texttt{/auth/me}；viewer 角色用户调用数据源管理接口。 & 无 token 返回 401；越权返回 403；不产生写入副作用。 \\
TC-03 & 发布审核状态流转（FR-03） & publisher 创建 draft 草稿→提交审核；admin 批准/驳回；重复提交。 & 状态按 draft→pending\_review→published/rejected 迁移；非法迁移返回 400；审核意见和审计日志存在。 \\
TC-04 & 报名/日历幂等（FR-04） & 同一用户对同一活动重复报名；重复加入日历；取消报名。 & 仅存在一条关联记录；取消后报名和日历事件同步移除。 \\
TC-05 & 数据源校验（FR-07） & 缺名称/URL 创建数据源；viewer 创建数据源；ftp 协议 URL。 & 参数错误 400；越权 403；拒绝非 http/https 协议和内网地址。 \\
TC-06 & 内容安全（FR-07） & 以 \texttt{<script>} 标签作为活动标题提交；抓取结果含控制字符。 & 输出已 HTML 转义；不执行脚本；控制字符被过滤。 \\
TC-07 & 检索与降级（FR-01） & 内部搜索空关键词、存在结果；外部搜索服务不可达。 & 内部返回稳定字段和排序；外部搜索降级返回 query/results/count/source/error 结构。 \\
TC-08 & 前端用户旅程 & 访客浏览活动列表→进入详情；登录后创建活动→提交审核；管理员审核。 & URL、视图、权限导航和错误态符合设计；E2E 可重复执行。 \\
TC-09 & 发布者申请（FR-05） & 普通用户提交发布者申请；管理员审批通过/驳回。 & 申请记录创建成功；审批后用户角色变更。 \\
TC-10 & 知识订阅通知（FR-06） & 用户订阅知识节点；管理员发布关联该节点的活动。 & 匹配活动发布后生成站内通知；用户可查看并标记已读。 \\
\bottomrule
\end{longtable}

\subsubsection{性能测试用例}
\begin{longtable}{p{1.0cm}p{2.8cm}p{4.5cm}p{4.5cm}}
\toprule 编号 & 场景 & 测试方法 & 预期指标 \\\midrule\endhead
TC-P01 & 活动列表响应时间 & 使用 API 测试工具连续请求活动列表 50 次，统计 P95 延迟。 & P95 < 500ms（本地/CI 环境）。 \\
TC-P02 & 登录接口响应时间 & 连续执行登录请求 30 次（含验证码校验），统计 P95 延迟。 & P95 < 1s（含验证码发送）。 \\
TC-P03 & 前端首页加载 & 使用 Playwright 测量首页首屏加载时间。 & 模拟校园网环境下 < 2s。 \\
TC-P04 & 并发请求稳定性 & 对活动列表接口发起 50 并发请求。 & 错误率 < 1\%。 \\
\bottomrule
\end{longtable}

\subsubsection{兼容性测试用例}
\begin{longtable}{p{1.0cm}p{2.8cm}p{4.5cm}p{4.5cm}}
\toprule 编号 & 场景 & 测试方法 & 预期结果 \\\midrule\endhead
TC-C01 & 视口宽度兼容 & 在 375px、768px、1024px、1920px 宽度下检查主要页面。 & 无水平溢出、无按钮遮挡、关键信息完整可见。 \\
TC-C02 & 浏览器兼容 & 在 Chromium（Playwright）中执行端到端测试。 & 全部 E2E 测试通过。 \\
TC-C03 & 键盘导航 & 使用 Tab 键遍历活动详情页、登录表单和管理页面。 & 焦点顺序合理，主要交互元素可键盘触达。 \\
\bottomrule
\end{longtable}

\subsubsection{异常与边界测试用例}
\begin{longtable}{p{1.0cm}p{2.8cm}p{4.5cm}p{4.5cm}}
\toprule 编号 & 场景 & 测试方法 & 预期结果 \\\midrule\endhead
TC-E01 & 网络异常 & 前端在网络断开状态下操作。 & 页面显示网络错误提示，不产生后台请求。 \\
TC-E02 & 空数据状态 & 无已发布活动时访问活动列表；无通知时访问通知页面。 & 显示空态提示和引导（如"暂无活动"）。 \\
TC-E03 & 非法输入 & 超长标题（>200 字符）；负数的活动 ID；非法日期格式。 & 服务端返回 400 和明确错误信息；前端显示校验提示。 \\
TC-E04 & SSRF 防护 & 提交 \texttt{file://} 协议、\texttt{10.0.0.1} 内网地址、\texttt{localhost} 的抓取目标 URL。 & 请求被拒绝，返回安全校验错误。 \\
TC-E05 & 并发重复提交 & 对同一活动快速连续发送两次报名请求。 & 数据库唯一约束保证仅一条报名记录。 \\
\bottomrule
\end{longtable}

\subsection{测试执行结果}
本次复核执行记录如下（2026-07-18）：

\noindent\begin{tabularx}{\textwidth}{p{2.5cm}Y p{4cm}}
\toprule 检查项 & 实际命令/观察 & 结论与后续动作 \\\midrule
前端聚合验证 & \texttt{npm run verify} 输出至 \texttt{vue-tsc -b \&\& vite build} 阶段，后续 unit/E2E 结果未出现 & 未完成。需在开发机或 CI 采集完整日志。 \\
前端单独构建 & \texttt{npm run build} 同样仅输出至 \texttt{vue-tsc -b \&\& vite build} & 构建阻塞待确认，需定位类型检查阻塞点。 \\
后端 pytest & \texttt{pytest -q -p no:cacheprovider} 返回 \texttt{No module named pytest} & 环境缺陷，需安装开发依赖后重跑。 \\
LaTeX 编译 & 六份报告 XeLaTeX 编译通过，图表经抽样检查 & 文档可正常生成。 \\
\bottomrule
\end{tabularx}

\section{缺陷管理}
\subsection{缺陷分类与严重度}
缺陷按严重度从高到低分为四个级别，每级附带具体示例和处置时限：
\begin{enumerate}
\item \textbf{P0（致命）：} 安全泄露（如 JWT 密钥或数据库密码提交到仓库）、数据破坏（如审核通过后活动内容丢失）、全站不可用（如 API 启动崩溃、数据库连接失败）。\textbf{处置：} 立即停止发布并回滚版本，24 小时内修复，修复后补充安全审查记录。
\item \textbf{P1（严重）：} 核心业务流程受阻且无可用替代方案。例如：登录接口返回 500 导致所有用户无法登录、活动提交审核后状态未更新（发布者无法确认审核进度）、报名接口报错（用户无法参与活动）。\textbf{处置：} 当前迭代内修复，必要时中断非核心功能开发。
\item \textbf{P2（一般）：} 非核心功能缺陷，存在可接受的替代方案。例如：搜索结果的排序不够精确、通知未实时推送但在下次登录时可同步、管理后台的筛选条件重置。\textbf{处置：} 可安排至下一迭代修复。
\item \textbf{P3（建议）：} 界面样式偏差、文案错误、文档过期。例如：按钮颜色与设计稿不一致、提示文案有错别字、API 文档中的示例字段与实际返回不匹配。\textbf{处置：} 视资源情况安排修复。
\end{enumerate}
缺陷按发现阶段分为编码阶段自测发现、代码审查发现、单元测试发现、集成测试发现、端到端测试发现和手工验收发现。越早发现缺陷修复成本越低，因此编码阶段的类型检查和代码审查是性价比最高的质量活动。

\subsection{缺陷生命周期与字段}
缺陷从发现到关闭经历以下状态流转：
\begin{center}
\begin{tikzpicture}[font=\small,
  state/.style={rectangle,rounded corners,draw,minimum width=2.2cm,minimum height=0.6cm,align=center,fill=blue!6},
  arr/.style={-{Stealth},thick}
]
\node[state] (new)       {新建};
\node[state] (confirmed) [right=1.5cm of new]      {已确认};
\node[state] (fixing)    [right=1.5cm of confirmed] {处理中};
\node[state] (verify)    [below=2.0cm of fixing]    {待验证};
\node[state] (reopen)    [below=0.8cm of confirmed]     {重开};
\node[state] (closed)    [below=1.2cm of reopen]     {已关闭};
\draw[arr] (new)--(confirmed);
\draw[arr] (confirmed)--(fixing);
\draw[arr] (fixing)--(verify);
\draw[arr] (verify)--node[above]{通过}(closed);
\draw[arr] (verify)--node[right]{不通过}(reopen);
\draw[arr] (reopen) -- (fixing);
\end{tikzpicture}
\vspace{0.2em}
{\par\small\itshape 图2：缺陷状态流转}
\end{center}

每个缺陷单记录以下完整字段：
\begin{enumerate}
\item \textbf{基本信息：} 缺陷编号（格式 QA-xxx）、发现版本号、发现环境（开发/测试/预发布）、发现阶段（编码/审查/单测/集成/E2E/验收）、报告人和报告日期。
\item \textbf{关联追踪：} 关联需求编号（FR-xx）或关联测试用例编号（TC-xx），用于需求—测试—缺陷的双向追溯。例如缺陷"登录验证码过期时间不符"关联 FR-02 和 TC-01。
\item \textbf{复现信息：} 复现步骤（Given-When-Then 格式）、实际结果、期望结果、日志片段或截图附件。
\item \textbf{分级与指派：} 严重度（P0--P3）、优先级（高/中/低）、负责人、目标修复版本。
\item \textbf{验证与关闭：} 修复验证人、验证日期、关联的回归测试名称或 ID。已关闭缺陷必须关联回归测试用例，或明确说明不适合自动化验证的原因。
\end{enumerate}

\section{质量保证方法}
\subsection{代码审查}
每个合并请求至少由一名非作者成员审查，使用以下逐项检查清单：
\begin{enumerate}
\item \textbf{权限安全：} 受限接口是否在服务端通过 JWT + 角色装饰器校验？前端是否根据角色隐藏或禁用越权入口？新增接口是否默认鉴权而非默认开放？
\item \textbf{异常路径：} 代码是否覆盖以下异常场景——空数据（列表/详情返回空态）、网络失败（请求超时或断开显示错误提示）、参数非法（后端返回 400 和明确错误信息）、越权访问（返回 403 且不产生写入副作用）？
\item \textbf{代码质量：} 命名是否清晰表达意图？函数是否超过 50 行需要拆分？条件逻辑嵌套是否超过三层需要提前返回？注释是否解释"为什么"而非"是什么"？
\item \textbf{测试覆盖：} 新增功能是否有对应的单元测试？涉及跨端流程是否有端到端测试？测试是否覆盖正向路径和至少一个异常路径？是否有被 `.skip` 跳过但未说明原因的测试用例？
\item \textbf{配置安全：} 是否引入了新的敏感信息（API Key、密码）？是否有硬编码的环境相关值（数据库地址、域名）？\texttt{.env.example} 是否同步更新？
\end{enumerate}
审查结果记录在合并请求评论中，所有检查项确认通过后方可合并。CI 失败、含真实密钥或破坏主干构建的合并请求禁止合并。紧急缺陷修复可走快速通道合并，但必须在修复后 24 小时内补齐测试用例和审查记录。

\subsection{静态检查}
前端在构建时通过 \texttt{vue-tsc -b} 执行 TypeScript 严格模式类型检查（\texttt{tsconfig.json} 中 \texttt{strict: true}），自动捕获以下问题：
\begin{enumerate}
\item 接口字段不匹配：前端 API 响应类型定义与组件期望的属性类型不一致时编译报错。
\item 空值引用：变量可能为 \texttt{null} 或 \texttt{undefined} 时未经判空直接访问其属性。
\item 类型错误：函数参数类型不符、枚举值超出定义范围、JSON 解析结果未指定类型。
\end{enumerate}
后端计划集成 pylint 或 flake8 在 CI 中执行静态检查，当前阶段以人工审查补充。静态检查门禁的目标是在编码阶段捕获 80\% 以上的类型和接口相关问题，降低运行时缺陷。

\subsection{自动化测试}
三项自动化测试手段已在项目中落地，通过 \texttt{npm run verify} 聚合执行：
\begin{enumerate}
\item \textbf{Vitest 单元测试（5 个测试文件）：} 覆盖认证 Store 的登录/登出状态变化、活动排序函数的多条件排序逻辑、活动质量评分函数的评分计算规则、日历卡片组件的渲染和事件绑定、活动详情组件的加载/就绪/错误三种状态分支。
\item \textbf{Playwright 端到端测试（1 个测试文件）：} 使用 Chromium 浏览器在 Mock 服务模式下验证访客浏览活动列表、用户注册和登录流程、受保护页面的未登录跳转。测试用例覆盖页面跳转、权限导航、表单校验提示和错误态展示。
\item \textbf{pytest 集成测试（20+ 测试文件）：} 使用 Flask 测试客户端和内存 SQLite，覆盖 auth（登录/注册/用户信息）、posters（活动 CRUD/状态流转）、audit（日志写入和查询）、data sources（创建/校验/列表）、knowledge（节点管理/关联）、search（内部搜索/外部降级）、quality（质量评分）、dedup（疑似重复检测）、security（URL 校验/HTML 转义）等模块。
\end{enumerate}

\subsection{回归测试}
自动化回归测试在以下三种场景中自动触发：
\begin{enumerate}
\item \textbf{每次 push：} CI 执行全部单元测试和端到端测试，确保提交不破坏已有功能。
\item \textbf{Pull request 创建/更新：} 在 push 触发的基础上，合并前必须全部通过。
\item \textbf{缺陷修复验证：} 关联缺陷的回归测试用例在修复提交时自动执行，确认修复有效且不引入新问题。
\end{enumerate}
端到端测试覆盖访客浏览、登录注册和活动搜索等关键用户旅程，确保功能迭代不破坏已有正确行为。回归测试的通过率为合并请求的硬性门禁——未全部通过不得合并。

\subsection{发布检查}
版本发布前的质量门禁按以下检查清单逐项确认：

\noindent\begin{tabularx}{\textwidth}{p{4cm}Y p{4cm}}
\toprule 检查项 & 确认方式 & 门禁标准 \\\midrule
CI 工作流 & GitHub Actions 运行结果 & 全部通过（构建 + 单测 + E2E） \\
缺陷状态 & 缺陷跟踪记录确认 & 目标版本 P0/P1 为零 \\
接口契约 & 契约矩阵与实际 API 对比 & 接口变更后已更新矩阵文档 \\
部署冒烟 & Docker Compose 启动 + 健康检查 & \texttt{/api/health} 返回 200 \\
版本标识 & Git tag + release notes & 版本 tag 已创建，release notes 已编写 \\
\bottomrule
\end{tabularx}
所有检查项确认通过后，由负责人在 release notes 中签署发布确认。任一检查项未达标则暂停发布，修复达标后方可继续。

\section{测试结论}
通过上述测试策略和质量保证方法，逸仙活动云项目建立了覆盖单元、集成、系统和验收四个层次的测试体系：
\begin{enumerate}
\item \textbf{测试资产完备：} 前端已配置 5 个 Vitest 单元测试文件和 1 个 Playwright 端到端测试文件；后端已编写 20+ pytest 测试文件覆盖认证、活动、审计、数据源、知识、搜索、质量、去重和安全等全部核心模块。
\item \textbf{自动化门禁已建立：} 前端 CI 工作流在 push 和 PR 时自动执行构建、类型检查、单元测试和端到端测试；\texttt{npm run verify} 提供提交前的一键聚合验证入口。
\item \textbf{需求-测试可追溯：} FR-02 对应认证 API 测试、auth store 测试和 \texttt{test\_auth\_api.py}；FR-03 对应活动编辑/审核页、posters API 测试和审计测试；FR-05/FR-07 对应知识、去重、质量、搜索和安全服务测试；FR-08 对应 Axios 拦截器、页面状态和 Playwright 授权流程。
\end{enumerate}
当前限制为执行环境的配置问题（QA-001：前端验证命令在本次检查中未完成；QA-002：后端环境缺少 pytest 依赖），这两项均为依赖安装问题而非测试代码缺陷，安装对应依赖后即可执行。性能基线测试、浏览器无障碍扫描和契约自动化测试列为后续补充项，不影响当前核心功能验证的可重复性。

\end{document}
